\documentclass[11pt]{article}
\usepackage[margin=0.82in]{geometry}
\usepackage{amsmath,amssymb}
\usepackage[colorlinks=true,linkcolor=blue,urlcolor=blue,citecolor=blue]{hyperref}
\setlength{\parindent}{0pt}
\setlength{\parskip}{0.62em}

\newcommand{\R}{\mathbb R}
\newcommand{\C}{\mathbb C}
\newcommand{\Z}{\mathbb Z}
\newcommand{\N}{\mathbb N}
\newcommand{\T}{\mathbb T}

\title{Irregular sampling overcomes the STFT lattice barrier\\
\large Exact affirmative answer in later literature}
\author{Run \texttt{fa\_banach\_001} --- \texttt{agent\_lane\_03}}
\date{11 August 2026}

\begin{document}
\maketitle

\textbf{Status.} \texttt{literature\_already\_answered}.  This packet
identifies and checks an exact later-literature answer; it does not claim a
new theorem.

\section*{Source question}

For a window $g\in L^2(\R^d)$, write
\[
 V_gf(x,\omega)=\int_{\R^d}f(t)\overline{g(t-x)}
 e^{-2\pi i\omega\cdot t}\,dt,
 \qquad f\sim h\iff f=\tau h\ \text{for some }\tau\in\T.
\]
Grohs and Liehr, \emph{Non-uniqueness theory in sampled STFT phase
retrieval}, arXiv:2207.05628v2, prove that for every window and every
ordinary lattice there are $f\not\sim h$ with identical sampled
spectrograms.  Their Conclusion (PDF page 32) then asks whether irregular
sampling, or a suitable increase in sampling redundancy, can restore
uniqueness.

The two ``natural question'' signals in the Introduction are not open: the
paper answers those lattice questions negatively.  The final irregular-
sampling question is the genuine forward-looking target.

\section*{Exact later answer}

Two months after the source preprint, the same authors posted
\emph{Phaseless sampling on square-root lattices}, arXiv:2209.11127.  Put
\[
 \sqrt{\Z}=\{\pm\sqrt n:n\in\N_0\}.
\]
Its Theorem 1.1 (PDF page 3) proves that a broad, dense class of analytic
windows admits phase retrieval of every $f\in L^2(\R^d)$ from a sufficiently
dense rectangular square-root lattice
\[
 \Lambda=A(\sqrt{\Z})^{2d}.
\]
Thus $|V_gf(\lambda)|=|V_gh(\lambda)|$ for all $\lambda\in\Lambda$ implies
$f\sim h$.  The set is discrete and irregular: coordinate gaps shrink away
from the origin, so it is not an ordinary lattice.

An explicit Gaussian specialization follows from Corollary 1.2.  If
$\varphi(x)=e^{-\gamma\lVert x\rVert_2^2}$ and
\[
 0<\alpha<\sqrt{\frac1{2\gamma e}},\qquad
 0<\beta<\sqrt{\frac{\gamma}{2\pi^2e}},
\]
then
\[
 |V_\varphi f|=|V_\varphi h|
 \quad\text{on}\quad
 \alpha(\sqrt\Z)^d\times\beta(\sqrt\Z)^d
 \qquad\Longleftrightarrow\qquad f\sim h.
\]
Corollary 1.3 also covers every Hermite window.

\section*{Why the answer is exact}

The source asks whether irregular sampling can be beneficial after proving
that no ordinary lattice can work.  The later theorem exhibits an explicit
irregular discrete family that does work, for unrestricted signals in
$L^2(\R^d)$.  It also cites the source paper as the lattice obstruction being
overcome.  Therefore the existential question has a full affirmative answer.

The answer is not a classification: it does not characterize all sampling
sets or all windows, and it does not say that merely increasing the density
of an ordinary lattice helps.  Indeed the source rules out that latter
possibility for every density.

\section*{Proof mechanism and intuition}

For the allowed analytic windows, the later paper proves that
$|V_gf|^2$ extends from $\R^{2d}$ to an entire function with explicit
coordinatewise Gaussian growth.  Jensen-type zero counting shows that a
sufficiently dense square-root lattice is a uniqueness set for precisely
that entire-function class.  Equality of the sampled magnitudes therefore
forces equality of the full spectrograms; a standard tensor/ambiguity
identity then forces $f\sim h$.

The key geometric point is that an ordinary lattice has constant spacing,
whereas the points $\pm\sqrt n$ become progressively denser.  Entire
functions with the prescribed growth cannot vanish on that many points
unless they vanish identically.  Irregularity supplies the information that
constant-density lattice sampling fundamentally lacks.

\section*{References and provenance}

[1] P. Grohs and L. Liehr, \emph{Non-uniqueness theory in sampled STFT
phase retrieval}, arXiv:2207.05628v2; \emph{SIAM J. Math. Anal.} 55 (2023),
4695--4726.  Question on PDF page 32.

[2] P. Grohs and L. Liehr, \emph{Phaseless sampling on square-root
lattices}, arXiv:2209.11127v2; \emph{Found. Comput. Math.} 25 (2025),
351--374.  Theorem 1.1 and Corollaries 1.2--1.3 on PDF pages 3--5.

The bundled PDFs are the official arXiv renderings.  Their SHA-256 values are
\texttt{574fc1fb...b986f9} and \texttt{69d314bf...7f84b}, respectively.

\end{document}
